\scanpage{76}
\noindent Soit donc \(U_i\) un recouvrement affine de \(X_0\), et
\[
 e\in H^{0^{(\mathrm{alg})}}(X_0,\mathscr{E}xt^1(K,\mathscr{N}))
\]
représenté par des classes d'isomorphies (algébriques) d'extensions \(\inserted{U'_i}\) de \(U_i\) par \(\mathscr{N}\mid U_i\), classes d'isomorphies qui se recollent :

\[
 \text{soit }\unsure{\eta|n|\gamma}_{ij}\text{ un isomorphisme}\qquad
 U'_{\unsure{j|i}}\mid U_i\cap U_j
 \xrightarrow{\sim}
 U'_{\unsure{i|j}}\mid U_i\cap U_j\ .
\]

L'obstruction à construire des \(\eta_{ij}\) \deleted{\unsure{d}}tels que \(\eta_{ij}\eta_{jk}=\eta_{ik}\) est \unsure{donné|donnée} par \(d_2\eta\) où \(d_2\) est une certaine flèche
\[
 d_2:H^0(X_0,\mathscr{E}xt^1(K,\mathscr{N}))\longrightarrow
 H^2(X_0,\mathscr{H}om(K,\mathscr{N}))
\]

La même construction marche dans le cas analytique, et on conclut par une nouvelle application de GAGA.

\uline{Rmq} \unsure{1|s} : \unsure{itou} pour les espaces algébriques, espaces rigides analytiques ; schémas formels

\deleted{\uline{\unsure{Rmq 2|Rmqs}}} :

\noindent\hfill Bien à toi\\
\null\hfill P. Deligne\\
\margin{\unsure{:}}
% archive sheet number: 75
\clearpage
